\documentclass[11pt, a4paper]{article}

\usepackage[T1]{fontenc}
\usepackage[utf8]{inputenc}
\usepackage{tgpagella}
\usepackage{tgcursor}
\usepackage{microtype}
\usepackage[spanish, es-noshorthands]{babel}
\addto\captionsspanish{\renewcommand{\tablename}{Tabla}}

\usepackage{amsmath, amssymb}
\usepackage[table, xcdraw]{xcolor}
\usepackage{booktabs}
\usepackage[most]{tcolorbox}
\usepackage[margin=2.5cm]{geometry}
\usepackage{fancyhdr}

\pagestyle{fancy}
\fancyhf{}
\setlength{\headheight}{16pt}
\lhead{\textbf{UCB Sede Tarija} | Ing. Mecatrónica}
\chead{Solucionario EC1 - Versión B}
\rhead{Robótica (IMT-342)}
\lfoot{Prof: M.Sc. Ing. Bernardo Quiroga Turdera}
\rfoot{Página \thepage}
\renewcommand{\headrulewidth}{0.4pt}
\renewcommand{\footrulewidth}{0.4pt}

\definecolor{ucbblue}{RGB}{0, 51, 102}

\begin{document}

\begin{center}
    \Large\textbf{\textcolor{red!80!black}{SOLUCIONARIO Y GUÍA DE CALIFICACIÓN (EC1 Versión B)}}\\[0.2cm]
    \normalsize Documento Confidencial Docente[cite: 8]
\end{center}

\section*{Problema 1 (20 pts)[cite: 8]}
\textbf{a) [6 pts]} ${}^{0}T_2 = {}^{0}T_1 {}^{1}T_2 \implies {}^{0}T_1 = {}^{0}T_2 ({}^{1}T_2)^{-1}$. (\textit{3 pts cadena correcta, 3 pts despeje}). \\
\textbf{b) [8 pts]} Inversa de ${}^{1}T_2$: Invertir $R$ (Identidad) y aplicar $-R^T p = -I[1,0,0]^T = [-1,0,0]^T$.
\begin{equation*}
    ({}^{1}T_2)^{-1} = \begin{bmatrix} 1&0&0&-1 \\ 0&1&0&0 \\ 0&0&1&0 \\ 0&0&0&1 \end{bmatrix} \implies
    {}^{0}T_1 = \begin{bmatrix} 0&-1&0&1 \\ 1&0&0&2 \\ 0&0&1&0 \\ 0&0&0&1 \end{bmatrix} \begin{bmatrix} 1&0&0&-1 \\ 0&1&0&0 \\ 0&0&1&0 \\ 0&0&0&1 \end{bmatrix} = \begin{bmatrix} 0&-1&0&1 \\ 1&0&0&1 \\ 0&0&1&0 \\ 0&0&0&1 \end{bmatrix}
\end{equation*}
(\textit{3 pts inversa, 3 pts multiplicacion, 2 pts resultado}).\\
\textbf{c) [6 pts]} El vector de traslación de ${}^{1}T_2$ está expresado en el marco 1. Antes de sumarlo con coordenadas del marco 0, debe ser rotado según la orientación del marco 1 respecto al 0. (\textit{Reconoce marco distinto 3 pts; rol de rotación 2 pts}).

\section*{Problema 2 (25 pts)[cite: 8]}
\textbf{a) [7 pts]} 1) El desplazamiento $a=2$ no es constante, requiere rotación $\cos q$ y $\sin q$. 2) El desplazamiento $d=1$ va en el componente Z, no Y. (\textit{Otorgar 7 pts por 2 errores claramente explicados}).\\
\textbf{b) [10 pts]} Para $\alpha=0, a=2, d=1$:
\begin{equation*}
    {}^{i-1}T_i = \begin{bmatrix} \cos q_i & -\sin q_i & 0 & 2\cos q_i \\ \sin q_i & \cos q_i & 0 & 2\sin q_i \\ 0 & 0 & 1 & 1 \\ 0 & 0 & 0 & 1 \end{bmatrix}
\end{equation*}
(\textit{3 pts $R$, 3 pts traslación $a$, 2 pts traslación $d$, 2 pts última fila}).\\
\textbf{c) [4 pts]} En $q_i = \pi/2$: ${}^{i-1}T_i = \begin{bmatrix} 0 & -1 & 0 & 0 \\ 1 & 0 & 0 & 2 \\ 0 & 0 & 1 & 1 \\ 0 & 0 & 0 & 1 \end{bmatrix}$.\\
\textbf{d) [4 pts]} $q_i = \theta_i$ porque es una articulación revoluta.

\section*{Problema 3 (45 pts)[cite: 8]}
\textbf{a) [8 pts]} Para $q_3=\pi, d_3=0, a_3=1, \alpha_3=0$: ${}^{2}T_3 = \begin{bmatrix} -1&0&0&-1 \\ 0&-1&0&0 \\ 0&0&1&0 \\ 0&0&0&1 \end{bmatrix}$. (\textit{3 pts R, 3 pts p, 2 pts fila base}). \\
\textbf{b) [5 pts]} ${}^{0}T_3 = {}^{0}T_1 {}^{1}T_2 {}^{2}T_3$.\\
\textbf{c) [16 pts]} ${}^{0}T_2 = {}^{0}T_1 {}^{1}T_2 = \begin{bmatrix} 0&0&1&0 \\ 1&0&0&2 \\ 0&1&0&0 \\ 0&0&0&1 \end{bmatrix}$. Luego ${}^{0}T_3 = {}^{0}T_2 {}^{2}T_3 = \begin{bmatrix} 0&0&1&0 \\ -1&0&0&1 \\ 0&-1&0&0 \\ 0&0&0&1 \end{bmatrix}$. (\textit{6 pts matriz ${}^0T_2$, 6 pts rotacion final, 4 pts posicion final}).\\
\textbf{d) [8 pts]} ${}^{0}p_3 = [0, 1, 0]^T$ (Origen marco 3 en base 0). ${}^{0}R_3 = \begin{bmatrix} 0&0&1 \\ -1&0&0 \\ 0&-1&0 \end{bmatrix}$ (Orientación de ejes 3 respecto a 0).\\
\textbf{e) [8 pts]} Estructural: Fila 4 es $[0,0,0,1]$ y $\det(R)=1$. Física: Evaluar magnitud espacial respecto a la secuencia de sumas de traslación de la matriz.

\section*{Problema 4 (10 pts)[cite: 8]}
\textbf{a) [5 pts]} Falla estructural: La última fila de la matriz propuesta es $[0, 0, 1, 1]$, lo cual es inválido. Debe ser estrictamente $[0, 0, 0, 1]$ para una transformación homogénea de cuerpo rígido.\\
\textbf{b) [5 pts]} Una traslación plausible no garantiza que la submatriz de rotación sea ortonormal y válida. Verificación adicional: Confirmar que $R^T R = I$ y $\det(R) = 1$.

\end{document}
